
\section{Detailed Step-by-Step Coding Manual}

Follow these steps \textbf{in order} inside \texttt{session5/session5\_iris\_template.py}.

\begin{tcolorbox}[title=Do: Add Imports]
Add imports at the very top:
\begin{minted}[fontsize=\small, linenos, frame=lines]{python}
import csv
from helper.iris_utils import fit_threshold_from_setosa
\end{minted}
\end{tcolorbox}

\begin{tcolorbox}[title=Do: Update Data Loading to Support Defaults]
Change from Session 4 style to Session 5 style:
\begin{minted}[fontsize=\small, linenos, frame=lines]{python}
def load_iris_data(filepath=None):
    if filepath is None:
        print("[WARNING] No filepath provided. Using default.")
        filepath = "session5/iris_data.csv"
\end{minted}
\end{tcolorbox}

\begin{tcolorbox}[title=Do: Add Optional Output Control in initialize\_predictions]
Change the function header to:
\begin{minted}[fontsize=\small, linenos, frame=lines]{python}
def initialize_predictions(dataset, settings, print_each=True):
\end{minted}
Then wrap your print statement inside \texttt{if print\_each:} so the caller can switch logging on or off.
\end{tcolorbox}

\begin{tcolorbox}[title=Do: Upgrade setup\_application]
Use this Session 5 header:
\begin{minted}[fontsize=\small, linenos, frame=lines]{python}
def setup_application(filepath=None, threshold=None, print_each=True, use_fit=False):
\end{minted}
Apply priority in this order:
\begin{enumerate}
    \item if \texttt{use\_fit} is True, learn threshold from \texttt{fit\_threshold\_from\_setosa(...)};
    \item else if \texttt{threshold} is not None, use the manual threshold override;
    \item else keep default threshold from \texttt{get\_default\_settings()}.
\end{enumerate}
\end{tcolorbox}


\section{Task: Positional vs Keyword Arguments in \texttt{main()}}

\textbf{Positional Arguments:} Python matches arguments to parameters by left-to-right order.
\begin{minted}[fontsize=\small, linenos, frame=lines]{python}
# filepath, threshold, print_each
setup_application("session5/iris_data.csv", 1.8, False)
\end{minted}

\textbf{Keyword Arguments:} You explicitly name the parameter you are setting.
\begin{minted}[fontsize=\small, linenos, frame=lines]{python}
setup_application(threshold=2.2, print_each=False)
\end{minted}

\begin{tcolorbox}[title=Do: Write main() with Four Runs]
In your \texttt{main()} function, add a helper \texttt{print\_summary} to avoid duplicate print statements. Then, write four distinct runs:
\begin{enumerate}
    \item \textbf{Default Run:} \texttt{setup\_application()}
    \item \textbf{Positional Run:} \texttt{setup\_application("session5/iris\_data.csv", 1.8, False)}
    \item \textbf{Keyword Run:} \texttt{setup\_application(threshold=2.2, print\_each=False)}
    \item \textbf{Fit Run:} \texttt{setup\_application(print\_each=False, use\_fit=True)}
\end{enumerate}
\end{tcolorbox}

\begin{tcolorbox}[title=Checkpoint]
Run the script using \texttt{python3 session5/session5\_iris\_template.py}.
You should see 4 separate summaries printed. The default run will print all 50 sample trace lines, while the other runs will not (because \texttt{print\_each=False}). 
\end{tcolorbox}

\section{Troubleshooting Mini-Section}

\begin{tcolorbox}[title=Common Mistakes \& Errors]
\begin{itemize}
    \item \textbf{ModuleNotFoundError:} If Python complains it cannot find \texttt{helper}, it means you ran the script from inside the \texttt{session5} directory. Navigate to the project root and run it from there.
    \item \textbf{SyntaxError: positional argument follows keyword argument:} In Python, if you mix positional and keyword arguments, all positional arguments \textbf{must} come before keyword arguments.
    \item \textbf{TypeError: multiple values for argument:} This happens if you accidentally pass a positional argument that maps to a parameter, and then provide a keyword argument for the exact same parameter.
\end{itemize}
\end{tcolorbox}

\section{Review and Next Steps}

\textbf{Summary:}
\begin{itemize}
    \item \textbf{Positional Arguments:} Matched by order. Easy to type, easy to mix up.
    \item \textbf{Keyword Arguments:} Matched by name. Self-documenting and safe.
    \item \textbf{Defaults (\texttt{None}):} Allow functions to be called without specifying every parameter, making APIs friendly and backward-compatible.
\end{itemize}

\textbf{Review Questions:}
\begin{enumerate}
    \item How do positional and keyword arguments differ?
    \item What happens if you pass a keyword argument before a positional argument?
    \item Why is it common to use \texttt{None} as a default value instead of \texttt{0} or \texttt{""}?
\end{enumerate}

\textbf{Micro-Exercise:} Change the keyword run in \texttt{main()} to use \texttt{threshold=1.5} and see how the accuracy changes. Try swapping the order of the keyword arguments (e.g., \texttt{print\_each=False, threshold=1.5}); does it still work? (Hint: Yes, keyword arguments don't depend on order!)

\textbf{Stretch Exercise:} Add an optional parameter \texttt{custom\_fit\_function=None} to \texttt{setup\_application}. If provided, use it instead of \texttt{fit\_threshold\_from\_setosa}.
\section{Importing Packages and Functions}
In Python, as your project grows, you will want to split your code into multiple files. You can import functions from one file into another using the \texttt{import} statement.

For this session, we provide a helper module \texttt{session5/helper/iris\_utils.py} which contains a function to calculate a threshold from the dataset.
You will \textbf{not} rewrite this helper; you will \textbf{use} it.
Import it at the top of your script like this:

\begin{minted}[fontsize=\small, linenos, frame=lines]{python}
import csv
from helper.iris_utils import fit_threshold_from_setosa
\end{minted}

\noindent
\textbf{Why this import is important for beginners:}
\begin{itemize}
    \item In Session 4, all logic stayed in one file.
    \item In Session 5, we start separating logic across files so code is easier to maintain.
    \item \texttt{helper.iris\_utils} means: "go to the \texttt{helper} folder and use the function inside \texttt{iris\_utils.py}".
\end{itemize}

\begin{tcolorbox}[title=Important: Running from the Project Root]
Beginners often think the import path is the same as the file path.
In Python, the import path depends on \textit{where you run the script from}.
For this tutorial, run your script from the project root:
\texttt{python3 session5/session5\_iris\_template.py}. This makes the \texttt{helper} module importable because the script looks for \texttt{helper} relative to where it was executed.
\end{tcolorbox}


\subsubsection{Your Task (Optional): Refactoring Code for Flexibility Using Keyword Arguments}

As an extension of the last session topic about conditional statement (if-else), I have create another scenario where we can apply the if-else statement to determines the antenna color of a robot based on its center\_x coordinate and a threshold value. Here is \href{https://github.com/balandongiv/computer_programming_gist/blob/main/session5/session5_robot_template.py#L43}{one way} on how to solve it:

\begin{minted}[fontsize=\small, linenos, breaklines=true]{python}
def determine_antenna_color(i, colors, config, my_th):
    if config["center_x"] >= my_th:
        return "purple"
    else:
        return colors[f"robot{i + 1}"]["light_color"]
\end{minted}

In this part of the code, we initially set the value of \texttt{my\_th} to \href{https://github.com/balandongiv/computer_programming_gist/blob/main/session5/session5_robot_template.py#L135}{500}, and this value was defined within the \href{https://github.com/balandongiv/computer_programming_gist/blob/main/session5/session5_robot_template.py#L121}{function} as shown below:

\begin{minted}[fontsize=\small, linenos, breaklines=true]{python}
def initialize_robots(canvas, data):
    colors = data["colors"]
    robot_configurations = data["robot_configurations"]
    condition = True
    my_th = 500
    for i, config in enumerate(robot_configurations):
        create_robot(i, canvas, config, condition, colors, my_th)
\end{minted}
However, hardcoding values, such as \texttt{my\_th}, is generally not a best practice in the real-world programming. This is because if the threshold value needs to be changed, you would have to manually update the code each time. A more flexible approach is to pass \texttt{my\_th} as a parameter. This method allows you to easily adjust the threshold value without needing to alter the code directly.

There are several ways to implement this improvement. For clarity and to focus on understanding keyword arguments, let's make changes starting from the \texttt{main()} function. Here's a step-by-step guide on how to refactor the code:
